\documentclass[11pt]{article}
\usepackage[margin=1.25in]{geometry}
\usepackage{amsmath,amssymb}
\usepackage[hidelinks]{hyperref}
\pagestyle{empty}
\setlength{\parskip}{6pt}
\begin{document}

\noindent
Claes Johnson\\
KTH Royal Institute of Technology\\
SE-100 44 Stockholm, Sweden\\
\texttt{claesjohnson@gmail.com}\\[1.2em]
\today\\[1.6em]

\noindent Dear Professor Fefferman,\\[0.6em]

I write only to let you know of an article I have posted to arXiv, \emph{Constructive Solution of the
Clay Navier--Stokes Problem} (arXiv:XXXX.XXXXX), since it engages directly with your formulation of
the problem and quotes you in it.

The article presents what an automated adaptive finite element method delivers when directed at
alternative~(A): for given data at high Reynolds number, turbulent solutions exhibited by
construction, with the error in mean-value outputs such as drag and lift bounded by duality rather
than estimated. I am explicit that this is not a discharge of~(A) as you state it --- what is
obtained is constructive existence and qualitative uniqueness, not smoothness, the solutions in
question having gradients that grow like $\nu^{-1/2}$ while their velocities stay bounded.

I quote your sentence that the four alternatives are offered ``to give reasonable leeway to solvers
while retaining the heart of the problem'', and take the Institute's own description to locate that
heart in turbulence. My claim is simply that I have used the leeway to reach turbulence. The material
itself is not new: it was published in 2007, in a chapter of a book written with Johan Hoffman.

Sending this is courtesy rather than a submission, and nothing in it needs a reply.\\[1.2em]

\noindent Yours sincerely,\\[2em]

\noindent Claes Johnson

\end{document}
